% GENERATED FILE. Edit canonical lessons, then run npm run generate:books.
% canonical-source-hash: fnv1a64:f8225d46849276f2
% canonical-lessons: HI-C88-karunga, HI-C88-aaunga, HI-C88-aaega, HI-C88-aaenge, HI-C88-aaoge, HI-C88-baje, HI-C88-repaso-bhavishya, HI-C88-sintesis-bhavishya

\chapter{The Future, And The Two Letters It Needed}
\label{ch:hi-bhavishya}
\hlchaptermodality{96}

\begin{chapteropening}
\textbf{By the end of this chapter:} \emph{I can say what I will do, what somebody else will do, and arrange to meet at a named hour.}

Read reading item 8 of the first practice paper word for word -- main panch baje aunga -- and answer the same question as a conversation at the aap level.
\end{chapteropening}

\section[karūṅgā]{\hi{करूँगा} — the third tense, and the one with no helper word}

\label{lesson:HI-C88-karunga}



\begin{quote}
You can say what you do and what you did. Say the stem of \emph{to do}, and hold it — the next tense attaches to it rather than standing beside it.
\end{quote}

\subsection*{You'll want to know: \hi{करूँगा}}
\textbf{\hi{करूँगा}} (\emph{karūṅgā}) — \textquotedblleft{}I will do,\textquotedblright{} said by a man. \textbf{\hi{करूँगी}} (\emph{karūṅgī}) — the same, said by a woman.

The stem is \textbf{\hi{कर}-}. Everything after it is the tense.

\begin{cousinweb}
Three pieces, and each one has a job:

\noindent
\begin{tabularx}{\linewidth}{@{}>{\raggedright\arraybackslash}X>{\raggedright\arraybackslash}X@{}}
\toprule
\textbf{piece} & \textbf{what it carries} \\
\midrule
\hi{कर}- & the action \\
-\hi{ऊँ}- & who is acting — \emph{I} \\
-\hi{गा} / -\hi{गी} & the future itself, agreeing with the speaker \\
\bottomrule
\end{tabularx}

English needs a second word for this: \emph{I \textbf{will} do}. Hindi has none. The tense is welded onto the verb, the way the person is welded into \emph{hūṁ} and \emph{hai}.

You have met the \textbf{-\hi{गा}} family once before, in the farewell \emph{milenge}, where the plural \textbf{-\hi{एंगे}} was named as the future ending. That was one cell of a grid. This chapter is the grid.
\end{cousinweb}

\begin{grammarlens}[title={Grammar Lens: the piece that comes from the old subjunctive}]
\textbf{-\hi{गा}} is not a random syllable. It comes from Sanskrit \textbf{\hi{गत}} (\emph{gata}), \textquotedblleft{}gone\textquotedblright{} — the participle of \emph{to go}. A future in Hindi is built out of a word meaning \textbf{gone}, which is the same instinct behind English \emph{I am going to do it}.

\noindent
\begin{tabularx}{\linewidth}{@{}>{\raggedright\arraybackslash}X>{\raggedright\arraybackslash}X@{}}
\toprule
\textbf{speaker} & \textbf{form} \\
\midrule
a man & \hi{करूँगा} \\
a woman & \hi{करूँगी} \\
\bottomrule
\end{tabularx}

Only the last letter moves, and it moves the way every \textbf{-\hi{आ} / -\hi{ई}} pair in this language moves.
\end{grammarlens}

\subsection*{Your turn}
Say these aloud:

\begin{itemize}
\raggedright
  \item \textquotedblleft{}\hi{करूँगा}\textquotedblright{} — \emph{karūṅgā}, I will do
  \item the same as a woman would say it
  \item the three pieces, and what each one carries
  \item how many words English needs for this, and how many Hindi needs
\end{itemize}

\begin{tcolorbox}[breakable,colback=teal!4,colframe=teal!35!black,title={Before you move on}]
Which piece carries the future? (\textbf{-\hi{गा}}, agreeing as \textbf{-\hi{गी}}.) What did its Sanskrit ancestor mean? (\emph{Gone}.) How many separate words does Hindi use to say \emph{will}? (None — it is an ending.)
\end{tcolorbox}
\section[āūṅgā]{\hi{आऊँगा} — the same ending, and the letter that was waiting for it}

\label{lesson:HI-C88-aaunga}



\begin{quote}
Say \emph{I will do}. Now take the stem of \emph{to come} and try to put the same ending on it — and listen for what happens when two vowels meet.
\end{quote}

\subsection*{You'll want to know: \hi{आऊँगा}}
\textbf{\hi{आऊँगा}} (\emph{āūṅgā}) — \textquotedblleft{}I will come,\textquotedblright{} said by a man. \textbf{\hi{आऊँगी}} (\emph{āūṅgī}) — the same, said by a woman.

Nothing new has been added. The stem \textbf{\hi{आ}-} has taken the ending you learned one lesson ago.

\begin{cousinweb}
\noindent
\begin{tabularx}{\linewidth}{@{}>{\raggedright\arraybackslash}X>{\raggedright\arraybackslash}X@{}}
\toprule
\textbf{stem} & \textbf{first person, future} \\
\midrule
\hi{कर}- & \hi{करूँगा} \\
\hi{आ}- & \hi{आऊँगा} \\
\bottomrule
\end{tabularx}

The ending is identical. What changes is how it has to be \textbf{written}, and that is worth a moment.

\textbf{\hi{कर}-} ends in a consonant, so the long \emph{ū} rides underneath it as the small mark you have used for a long time. \textbf{\hi{आ}-} ends in a vowel, so there is nothing to ride on — and the \emph{ū} has to be written as a \textbf{full letter of its own}.
\end{cousinweb}

\begin{grammarlens}[title={Grammar Lens: the letter arrives where it was promised}]
Find \textbf{\hi{ऊ}} inside the word:

\begin{quote}
\hi{आ} \textbf{\hi{ऊँ}} \hi{गा}
\end{quote}

That is the letter you were told was coming — the one with the extra loop on the right — and this is the word it was waiting for. A vowel that opens a syllable needs the full letter, and the first-person future of every vowel-stem verb puts one there.

The same thing happens to \emph{to go}:

\noindent
\begin{tabularx}{\linewidth}{@{}>{\raggedright\arraybackslash}X>{\raggedright\arraybackslash}X@{}}
\toprule
\textbf{verb} & \textbf{I will} \\
\midrule
\hi{आना} & \hi{आऊँगा} \\
\hi{जाना} & \hi{जाऊँगा} \\
\bottomrule
\end{tabularx}

Both of these appear on a practice paper, in the first scored part.
\end{grammarlens}

\subsection*{Your turn}
\begin{itemize}
\raggedright
  \item \emph{Say it:} \textquotedblleft{}\hi{आऊँगा}\textquotedblright{} — \emph{āūṅgā}, I will come
  \item \emph{Say it:} the same as a woman would say it
  \item \emph{Look:} at \hi{आऊँगा} and put your finger on the letter \textbf{\hi{ऊ}}
  \item \emph{Say it:} why this word needs the full letter where \emph{karūṅgā} needs only a mark
\end{itemize}

\begin{tcolorbox}[breakable,colback=teal!4,colframe=teal!35!black,title={Before you move on}]
What decides whether the long \emph{ū} is a mark or a full letter? (Whether the stem ends in a consonant or a vowel.) Say \emph{I will go}. (\emph{Jāūṅgā} — \textbf{\hi{जाऊँगा}}.)
\end{tcolorbox}
\section[āegā]{\hi{आएगा} — the ending that tells you who, not when}

\label{lesson:HI-C88-aaega}



\begin{quote}
You can say \emph{I will come}. Point at somebody else and try to say it about them — and notice which part of the word has to move.
\end{quote}

\subsection*{You'll want to know: \hi{आएगा}}
\textbf{\hi{आएगा}} (\emph{āegā}) — \textquotedblleft{}he will come.\textquotedblright{} \textbf{\hi{आएगी}} (\emph{āegī}) — \textquotedblleft{}she will come.\textquotedblright{}

\textbf{\hi{वह} \hi{आएगा।}} — He will come. \textbf{\hi{वह} \hi{आएगी।}} — She will come.

\begin{cousinweb}
Set the two first persons beside the third:

\noindent
\begin{tabularx}{\linewidth}{@{}>{\raggedright\arraybackslash}X>{\raggedright\arraybackslash}X@{}}
\toprule
\textbf{who} & \textbf{ending} \\
\midrule
\hi{मैं} & -\hi{ऊँगा} \\
\hi{वह} & -\hi{एगा} \\
\bottomrule
\end{tabularx}

The \textbf{-\hi{गा}} has not moved. It is the piece in front of it that changed, from \textbf{-\hi{ऊँ}-} to \textbf{-\hi{ए}-}, and that piece is the one carrying the person.

So the future has the same division of labour the present had: \textbf{one slot for who, one slot for when.} Learning the tense is learning that the two slots are separate.
\end{cousinweb}

\begin{grammarlens}[title={Grammar Lens: the form a paper prints}]
This is the commonest future in reading material, because most sentences on a page are about somebody other than the reader:

\noindent
\begin{tabularx}{\linewidth}{@{}>{\raggedright\arraybackslash}X>{\raggedright\arraybackslash}X@{}}
\toprule
\textbf{sentence} & \textbf{meaning} \\
\midrule
\hi{वह} \hi{आएगा।} & He will come. \\
\hi{वह} \hi{जाएगी।} & She will go. \\
\bottomrule
\end{tabularx}

A bus, a train or a shop takes the same ending as a person, and the gender of the \textbf{noun} decides which one. A practice paper prints \emph{the bus will come} with \textbf{-\hi{एगी}}, because \hi{बस} is feminine.
\end{grammarlens}

\subsection*{Your turn}
Say these aloud:

\begin{itemize}
\raggedright
  \item \textquotedblleft{}\hi{वह} \hi{आएगा}\textquotedblright{} — he will come
  \item \textquotedblleft{}\hi{वह} \hi{आएगी}\textquotedblright{} — she will come
  \item which slot changed between \emph{āūṅgā} and \emph{āegā}, and which stayed
  \item \textquotedblleft{}\hi{वह} \hi{जाएगी}\textquotedblright{}
\end{itemize}

\begin{tcolorbox}[breakable,colback=teal!4,colframe=teal!35!black,title={Before you move on}]
Which piece of the future word says \emph{who}? (The one before \textbf{-\hi{गा}}.) Which piece says \emph{when}? (\textbf{-\hi{गा}} itself.) What ending does a feminine noun take? (\textbf{-\hi{एगी}}.)
\end{tcolorbox}
\section[āeṅge]{\hi{आएँगे} — the ending you have been saying since the first goodbye}

\label{lesson:HI-C88-aaenge}



\begin{quote}
Say the farewell that means \emph{see you again}. The last two syllables of it are about to become a rule.
\end{quote}

\subsection*{You'll want to know: \hi{आएँगे}}
\textbf{\hi{आएँगे}} (\emph{āeṅge}) — \textquotedblleft{}we will come,\textquotedblright{} \textquotedblleft{}they will come,\textquotedblright{} and \textquotedblleft{}you will come\textquotedblright{} to somebody you address as \textbf{\hi{आप}}.

\textbf{\hi{हम} \hi{आएँगे।}} — We will come. \textbf{\hi{आप} \hi{कब} \hi{आएँगे}?} — When will you come?

\begin{cousinweb}
You met this ending in your fourth chapter, inside \textbf{\hi{मिलेंगे}} — \emph{we will meet}. It was named there as the plural future, and it was true, but it was one word. Now it is a slot you can fill:

\noindent
\begin{tabularx}{\linewidth}{@{}>{\raggedright\arraybackslash}X>{\raggedright\arraybackslash}X@{}}
\toprule
\textbf{stem} & \textbf{plural future} \\
\midrule
\hi{मिल}- & \hi{मिलेंगे} \\
\hi{आ}- & \hi{आएँगे} \\
\hi{कर}- & \hi{करेंगे} \\
\bottomrule
\end{tabularx}

The farewell was never a lump. It was this ending, on a verb meaning \emph{to meet}.
\end{cousinweb}

\begin{grammarlens}[title={Grammar Lens: \hi{आप} takes the plural, again}]
\noindent
\begin{tabularx}{\linewidth}{@{}>{\raggedright\arraybackslash}X>{\raggedright\arraybackslash}X@{}}
\toprule
\textbf{subject} & \textbf{future} \\
\midrule
\hi{हम} & \hi{आएँगे} \\
\hi{वे} & \hi{आएँगे} \\
\hi{आप} & \hi{आएँगे} \\
\bottomrule
\end{tabularx}

\textbf{\hi{आप}} takes the plural ending even for one person, for the reason it has taken a plural everywhere else: the polite level in this language is built out of plural forms. That is now the fifth system in which one decision about how you address somebody settles the whole row.

For a group of women, the ending nasalises differently and becomes \textbf{\hi{आएँगी}} — worth recognising, because a practice paper asks \emph{\hi{आप} \hi{कब} \hi{आएँगी}?} of a woman.
\end{grammarlens}

\subsection*{Your turn}
Say these aloud:

\begin{itemize}
\raggedright
  \item \textquotedblleft{}\hi{हम} \hi{आएँगे}\textquotedblright{} — we will come
  \item \textquotedblleft{}\hi{आप} \hi{कब} \hi{आएँगे}?\textquotedblright{} — when will you come?
  \item the farewell, and then name the ending inside it
  \item \textquotedblleft{}\hi{हम} \hi{करेंगे}\textquotedblright{}
\end{itemize}

\begin{tcolorbox}[breakable,colback=teal!4,colframe=teal!35!black,title={Before you move on}]
Which ending does \textbf{\hi{आप}} take in the future? (The plural one.) Which familiar farewell is built from it? (\emph{Milenge}.) What was that farewell really made of? (A stem plus this ending.)
\end{tcolorbox}
\section[āoge]{\hi{आओगे} — the \hi{तुम} row, and a vowel that has shown up before}

\label{lesson:HI-C88-aaoge}



\begin{quote}
Say the command you would give a friend using \textbf{\hi{तुम}}. Hold on to the vowel at its end — the future is about to reuse it.
\end{quote}

\subsection*{You'll want to know: \hi{आओगे}}
\textbf{\hi{आओगे}} (\emph{āoge}) — \textquotedblleft{}you will come,\textquotedblright{} to somebody you address as \textbf{\hi{तुम}}. \textbf{\hi{करोगे}} (\emph{karoge}) — \textquotedblleft{}you will do,\textquotedblright{} at the same level.

\begin{cousinweb}
The \textbf{\hi{तुम}} level has a vowel of its own, and you have already used it:

\noindent
\begin{tabularx}{\linewidth}{@{}>{\raggedright\arraybackslash}X>{\raggedright\arraybackslash}X@{}}
\toprule
\textbf{where} & \textbf{form} \\
\midrule
the command & \hi{बोल}\textbf{\hi{ो}} \\
the future & \hi{आ}\textbf{\hi{ओ}}\hi{गे} \\
\bottomrule
\end{tabularx}

Both are \textbf{-\hi{ओ}}. That is not two facts to learn — it is one vowel that belongs to \textbf{\hi{तुम}} wherever it appears, and you have had it since the commands chapter.
\end{cousinweb}

\begin{grammarlens}[title={Grammar Lens: where the full letter shows up again}]
The same writing question as the first person, and the same answer:

\noindent
\begin{tabularx}{\linewidth}{@{}>{\raggedright\arraybackslash}X>{\raggedright\arraybackslash}X@{}}
\toprule
\textbf{stem} & \textbf{you will} \\
\midrule
\hi{कर}- & \hi{करोगे} \\
\hi{आ}- & \hi{आओगे} \\
\bottomrule
\end{tabularx}

\textbf{\hi{करोगे}} carries the \emph{o} as a mark on the consonant. \textbf{\hi{आओगे}} cannot, because \textbf{\hi{आ}-} is a vowel, so the \emph{o} is written as the full letter \textbf{\hi{ओ}} — the one built on \textbf{\hi{आ}} with an arc above it.

Two vowel-stem words in one chapter, two full letters, and the same rule producing both.
\end{grammarlens}

\subsection*{Your turn}
\begin{itemize}
\raggedright
  \item \emph{Say it:} \textquotedblleft{}\hi{तुम} \hi{आओगे}\textquotedblright{} — you will come
  \item \emph{Say it:} \textquotedblleft{}\hi{तुम} \hi{करोगे}\textquotedblright{}
  \item \emph{Look:} at \hi{आओगे} and put your finger on the letter \textbf{\hi{ओ}}
  \item \emph{Say it:} the same question at the \hi{आप} level instead
\end{itemize}

\begin{tcolorbox}[breakable,colback=teal!4,colframe=teal!35!black,title={Before you move on}]
Which vowel belongs to the \textbf{\hi{तुम}} level? (\textbf{-\hi{ओ}}.) Where else have you used it? (On the command.) Why does \emph{āoge} need the full letter and \emph{karoge} does not? (The stem ends in a vowel, so nothing can carry the mark.)
\end{tcolorbox}
\section[pā̃c baje]{\hi{पाँच} \hi{बजे} — reading the hour, and arranging to meet at it}

\label{lesson:HI-C88-baje}



\begin{quote}
You can say what the clock is doing — the verb that strikes. Now try to say \emph{at five}, and notice that telling the time and fixing a time are not the same job.
\end{quote}

\subsection*{You'll want to know: \hi{पाँच} \hi{बजे}}
\textbf{\hi{पाँच} \hi{बजे}} (\emph{pā̃c baje}) — \textquotedblleft{}at five o'clock.\textquotedblright{}

\textbf{\hi{मैं} \hi{पाँच} \hi{बजे} \hi{आऊँगा।}} — I will come at five.

\begin{cousinweb}
\textbf{\hi{बजे}} comes from the verb you already have — the one meaning \emph{to strike, to toll}, because clocks and towns once marked the hour by striking a bell.

What is new is the \textbf{shape} of it. You learned it as something the clock is doing right now. This form is not doing anything: it sits beside a number and turns it into a \textbf{time to meet at}.

\noindent
\begin{tabularx}{\linewidth}{@{}>{\raggedright\arraybackslash}X>{\raggedright\arraybackslash}X@{}}
\toprule
\textbf{what it does} & \textbf{the form} \\
\midrule
tells you the hour & the striking form you know \\
fixes a time to act & \textbf{\hi{बजे}} \\
\bottomrule
\end{tabularx}

There is no separate word for \emph{at}. The number plus \textbf{\hi{बजे}} is the whole of it.
\end{cousinweb}

\begin{grammarlens}[title={Grammar Lens: where it stands in the sentence}]
\noindent
\begin{tabularx}{\linewidth}{@{}>{\raggedright\arraybackslash}X>{\raggedright\arraybackslash}X@{}}
\toprule
\textbf{sentence} & \textbf{meaning} \\
\midrule
\hi{मैं} \hi{पाँच} \hi{बजे} \hi{आऊँगा।} & I will come at five. \\
\hi{वह} \hi{चार} \hi{बजे} \hi{आएगी।} & She will come at four. \\
\bottomrule
\end{tabularx}

The time phrase goes \textbf{after the person and before the verb}, in the same seat the place phrase took. Hindi keeps the verb at the end and fills the middle with whatever the sentence needs.

Both practice papers ask for exactly this pairing — a number, \textbf{\hi{बजे}}, and a future verb — in the first scored part.
\end{grammarlens}

\subsection*{Your turn}
Say these aloud:

\begin{itemize}
\raggedright
  \item \textquotedblleft{}\hi{पाँच} \hi{बजे}\textquotedblright{} — at five
  \item \textquotedblleft{}\hi{मैं} \hi{पाँच} \hi{बजे} \hi{आऊँगा}\textquotedblright{}
  \item the same sentence as a woman would say it
  \item which word Hindi uses for \emph{at}, and hear that there is none
\end{itemize}

\begin{tcolorbox}[breakable,colback=teal!4,colframe=teal!35!black,title={Before you move on}]
Which verb is \textbf{\hi{बजे}} from? (The one meaning \emph{to strike}.) How does Hindi say \emph{at} in front of a clock time? (It does not — the number plus \textbf{\hi{बजे}} carries it.) Where does the time phrase stand? (After the person, before the verb.)
\end{tcolorbox}
\section[bhaviṣya kī tālikā]{bhaviṣya kī tālikā — the future, on one page}

\label{lesson:HI-C88-repaso-bhavishya}



\begin{quote}
Name the ending for \emph{I} and the ending for \emph{he}, and say which piece of each one is the tense, before you look at the table.
\end{quote}

\begin{grammarlens}[title={Grammar Lens: the grid}]
\noindent
\begin{tabularx}{\linewidth}{@{}>{\raggedright\arraybackslash}X>{\raggedright\arraybackslash}X@{}}
\toprule
\textbf{subject} & \textbf{with \hi{आ}-} \\
\midrule
\hi{मैं} & \hi{आऊँगा} / \hi{आऊँगी} \\
\hi{तुम} & \hi{आओगे} \\
\hi{वह} & \hi{आएगा} / \hi{आएगी} \\
\hi{हम}, \hi{वे}, \hi{आप} & \hi{आएँगे} / \hi{आएँगी} \\
\bottomrule
\end{tabularx}

Four rows, one verb, and \textbf{-\hi{गा}} in every single cell. What moves is the vowel in front of it, and that vowel is the person.
\end{grammarlens}

\begin{grammarlens}[title={Grammar Lens: two slots, three tenses}]
\noindent
\begin{tabularx}{\linewidth}{@{}>{\raggedright\arraybackslash}X>{\raggedright\arraybackslash}X@{}}
\toprule
\textbf{tense} & \textbf{how it is built} \\
\midrule
now & stem + -\hi{ता} + \hi{हूँ} / \hi{है} \\
earlier & stem + -\hi{ता} + \hi{था} \\
later & stem + person vowel + -\hi{गा} \\
\bottomrule
\end{tabularx}

The first two share a shape: the habit is in \textbf{-\hi{ता}} and the time is in the word after it. The future is the odd one out — it has \textbf{no separate word at all}, and packs the person and the tense into one ending.
\end{grammarlens}

\begin{grammarlens}[title={Grammar Lens: where the full letters appear}]
\noindent
\begin{tabularx}{\linewidth}{@{}>{\raggedright\arraybackslash}X>{\raggedright\arraybackslash}X@{}}
\toprule
\textbf{stem ends in} & \textbf{how the vowel is written} \\
\midrule
a consonant (\hi{कर}-) & a mark on the letter \\
a vowel (\hi{आ}-, \hi{जा}-) & a full letter — \textbf{\hi{ऊ}} or \textbf{\hi{ओ}} \\
\bottomrule
\end{tabularx}

That is the whole reason two new letters were needed before this chapter could be written.
\end{grammarlens}

\subsection*{Your turn}
Say these aloud:

\begin{itemize}
\raggedright
  \item all four rows of the grid, top to bottom
  \item the same grid with \textbf{\hi{कर}-} instead of \textbf{\hi{आ}-}
  \item \textquotedblleft{}\hi{मैं} \hi{पाँच} \hi{बजे} \hi{आऊँगा}\textquotedblright{}
  \item which two forms need a full letter, and why
\end{itemize}

\begin{tcolorbox}[breakable,colback=teal!4,colframe=teal!35!black,title={Before you move on}]
Which piece is in every cell of the grid? (\textbf{-\hi{गा}}.) Which tense uses no separate word for time? (The future.) How does Hindi say \emph{at five}? (\emph{Pā̃c baje} — no word for \emph{at}.)
\end{tcolorbox}
\section[maiṁ pā̃c baje āūṅgā]{\hi{मैं} \hi{पाँच} \hi{बजे} \hi{आऊँगा} — the item the paper prints}

\label{lesson:HI-C88-sintesis-bhavishya}



\begin{quote}
Read this out loud before reading on.

\begin{quote}
\hi{मैं} \hi{पाँच} \hi{बजे} \hi{आऊँगा।}
\end{quote}
\end{quote}

\begin{grammarlens}[title={Grammar Lens: four words, four things you were taught}]
That line is reading item 8 of the first practice paper, word for word. Until this chapter two of its four words were unreadable.

\noindent
\begin{tabularx}{\linewidth}{@{}>{\raggedright\arraybackslash}X>{\raggedright\arraybackslash}X@{}}
\toprule
\textbf{word} & \textbf{what it does} \\
\midrule
\hi{मैं} & who \\
\hi{पाँच} \hi{बजे} & when \\
\hi{आऊँगा} & the verb, carrying the tense \\
\bottomrule
\end{tabularx}

The answer the paper wants is \emph{I will come at five}, and the two wrong answers offered are past and present. \textbf{The only thing telling them apart is the ending on the last word.}
\end{grammarlens}

\subsection*{The exchange}
The second paper asks the same thing as a conversation:

\begin{quote}
— \hi{आप} \hi{कब} \hi{आएँगे}? — \hi{चार} \hi{बजे।}
\end{quote}

\textbf{\hi{आएँगे}} because the question is to somebody addressed as \textbf{\hi{आप}}. The answer needs no verb at all: a number and \textbf{\hi{बजे}} is a complete reply, the way \emph{at four} is in English.

\subsection*{How to answer}
Say each of these out loud:

\begin{itemize}
\raggedright
  \item you will come at five
  \item you will go by train, not by bus
  \item she will come in ten minutes
\end{itemize}

The second and third are items 13 and 11 of the same papers. All three are the same shape: a person, a circumstance, and a verb whose ending does the work.

\subsection*{Your turn}
Say these aloud:

\begin{itemize}
\raggedright
  \item \textquotedblleft{}\hi{मैं} \hi{पाँच} \hi{बजे} \hi{आऊँगा}\textquotedblright{}
  \item the same sentence as a woman would say it
  \item \textquotedblleft{}\hi{आप} \hi{कब} \hi{आएँगे}?\textquotedblright{} and answer it with two words
  \item \textquotedblleft{}\hi{वह} \hi{आएगी}\textquotedblright{}
\end{itemize}

\begin{tcolorbox}[breakable,colback=teal!4,colframe=teal!35!black,title={Before you move on}]
Which single word in the paper's line tells you it is the future? (\textbf{\hi{आऊँगा}} — its ending.) How many words does a full answer to \emph{when} need? (Two — the number and \textbf{\hi{बजे}}.) Which ending does a question to \textbf{\hi{आप}} take? (The plural, \textbf{\hi{आएँगे}}.)
\end{tcolorbox}
